\documentclass[Screen16to9,17pt]{foils}
\usepackage{zencurity-slides}

\selectlanguage{english}

\usepackage{tabularx}

% Content reused from pentest/pentest-II-foredrag. Goals/Plan from subject
% plan (obj 7). FIXME = topics not in that deck.

\begin{document}

\mytitlepage{Web Attacks I -- Attacks}
{Saint Cyber -- Cyber Akademiet, 2026}

\slide{Goals for today}
\begin{list2}
\item Identify and explain common web-application weaknesses
\item Cover SQL injection, cross-site scripting and cross-site request forgery
\item Cover session handling, file-inclusion flaws and security misconfiguration
\item Relate them to the OWASP Top 10
\end{list2}

\slide{Plan for today}
\begin{list2}
\item Why web hacking -- and the OWASP Top 10 as a map
\item Injection: SQL injection
\item Client-side: XSS and CSRF
\item Sessions, file inclusion and misconfiguration
\end{list2}

\slide{Why talk about hacking web}

\hlkimage{7cm}{pawel-janiak-dxFi8Ea670E-unsplash.jpg}

\begin{list2}
\item Everything is web today
\item HTTPS/TLS used for data transport in most applications
\item Web APIs are used across the industry
\item Even \emph{desktop} and \emph{mobile} application are rarely more than a specialized browser\\ like \url{electronjs.org} -- Chromium and Node.js
\hfill {\footnotesize Photo by Pawel Janiak on Unsplash}
\item Also software security and web security has a HUGE overlap
\end{list2}

\slide{Web Application Security}

\hlkimage{6cm}{hoffman-web-application-security.jpg}
\emph{Web Application Security}, Andrew Hoffman, 2020, ISBN: 9781492053118 called WAS

\slide{OWASP Top Ten}

\hlkimage{16cm}{owasp.jpg}

\begin{quote}
The OWASP Top Ten provides a minimum standard for web application
security. The OWASP Top Ten represents a broad consensus about what
the most critical web application security flaws are.
\end{quote}

\begin{list1}
\item The Open Web Application Security Project (OWASP)
\item OWASP produces lists of the most common types of errors in web applications
\item \link{http://www.owasp.org}
\item Create Secure Software Development Lifecycle
\end{list1}

\slide{SQL injection}

SQL Injection FAQ from Chip Andrews 2001 \link{http://www.sqlsecurity.com}:

\begin{minted}[fontsize=\footnotesize]{shell}
Set myRecordset = myConnection.execute
("SELECT * FROM myTable
WHERE someText ='" & request.form("inputdata") & "'")
\end{minted}

Using input \verb+exec master..xp_cmdshell 'net user test testpass /ADD' --+

Will perform
\begin{minted}[fontsize=\footnotesize]{shell}
SELECT * FROM myTable
WHERE someText ='' exec master..xp_cmdshell 'net user test testpass /ADD'--
\end{minted}

Depending on the access rights this might have full authorization in the database!

\slide{Cross-site scripting}

\begin{list1}
\item When logged into a site, we have a session identifier -- being presented with each HTTP request
\item If an attacker can \emph{activate} javascript that access and re-send this to another site we have a cross-site scripting attack
\item We often check using some parameter with code, like this:
\begin{alltt}\footnotesize
<A HREF="http://example.com/comment.cgi?
mycomment=<SCRIPT>malicious code</SCRIPT>
">Click here</A>
\end{alltt}
\item If this code is returned, and active, as part of the returned page and HTML we have a cross-site scripting possibility
\item This code is executed in the user browser, with the same permissions as the user
\end{list1}

\slide{11. Cross-Site Request Forgery (CSRF)}

\begin{quote}
Sometimes we already know an API endpoint exists that would allow us to perform an operation we wish to perform, but we do not have access to that endpoint because it requires privileged access (e.g., an admin account).

In this chapter, we will be building Cross-Site Request Forgery (CSRF) exploits that result in an admin or privileged account performing an operation on our behalf rather than using a JavaScript code snippet.

CSRF attacks take advantage of the way browsers operate and the trust relationship between a website and the browser. By finding API calls that rely on this relationship to ensure security—but yield too much trust to the browser—we can craft links and forms that with a little bit of effort can cause a user to make requests on his or her own behalf—unknown to the user generating the request.
\end{quote}

\begin{list2}
  \item Often seen in small CPE routers in homes, if the user activates and evil link, their router might be reconfigured or taken over
\end{list2}

\slide{Session IDs}

\begin{list2}
\item Session IDs tie the user with the state on the server
\item Must be randomly assigned, otherwise an attacker can guess a valid ID
\item Common problems, time based or predictable in some way
\item Check code for generating IDs or measure - Phase Space Analysis
\end{list2}

\slide{File inclusion (LFI / RFI)}
% FIXME not in pentest-II -- add local/remote file inclusion (subject plan obj 7)
\begin{list2}
\item \textbf{FIXME:} not in pentest-II -- add local/remote file inclusion (subject plan obj 7)
\end{list2}

\slide{Security misconfiguration}
% FIXME add defaults / verbose errors / missing headers (subject plan obj 7)
\begin{list2}
\item \textbf{FIXME:} add defaults / verbose errors / missing headers (subject plan obj 7)
\end{list2}

\slide{Summary}
\begin{list2}
\item The main web weaknesses: injection, XSS, CSRF, sessions, file inclusion, misconfiguration
\item Next: the tools, labs and defences
\end{list2}

\end{document}
